% =============================================================================
% ISI Country Brief — Portugal
% External Supplier Concentration Profile
%
% ISI Paper Series — Country Brief
% International Sovereignty Institute — Research Unit
% March 2026
% =============================================================================
\documentclass[12pt,a4paper]{article}

% --- ISI Paper Series shared template ----------------------------------------
\usepackage{isi-series}

% --- PDF metadata ------------------------------------------------------------
\hypersetup{
  pdftitle={External Supplier Concentration Profile: Portugal — ISI EU-27 2024},
  pdfauthor={International Sovereignty Institute -- Research Unit},
  pdfcreator={International Sovereignty Institute},
  pdfsubject={ISI Paper Series -- Country Brief PT, EU-27, Vintage 2024},
  pdfkeywords={sovereignty index, EU-27, supplier concentration,
    portugal, country brief},
}

% ==============================================================================
\begin{document}
\hypersetup{pageanchor=false}

% ---------- Title Page with Metadata ------------------------------------------
\begin{titlepage}
\centering
\vspace*{2cm}
{\Large\textsc{ISI Paper Series --- Country Brief}}\\[1.2cm]
{\LARGE\bfseries External Supplier Concentration Profile:\\[0.3em]
Portugal}\\[1.5cm]
{\large International Sovereignty Index (ISI)\\[0.3em]
EU-27 --- Vintage 2024 --- Methodology v1.0}\\[1.2cm]
{\large International Sovereignty Institute --- Research Unit}\\[0.8cm]
{\large March 2026}

\vspace{0.8cm}
\noindent\rule{\textwidth}{0.4pt}
\vspace{0.4cm}

% --- Research Metadata --------------------------------------------------------
\begin{flushleft}
\noindent\textbf{Data basis:} ISI Paper~2 --- EU-27 Empirical Results 2024 (Methodology~v1.0, Vintage~2024).\\[4pt]
\noindent\textbf{Methodology:} ISI Paper~1 --- Methodological Foundations, Version~1.0.

\medskip

\noindent\textbf{JEL Codes:} F02, F15, F52, O30, Q43\\[4pt]
\noindent\textbf{Keywords:} sovereignty index; EU-27; supplier concentration; composite indicator; variance decomposition; defence dependence

\medskip

\noindent\textbf{Related publications in the ISI Paper Series:}
\begin{itemize}[nosep,leftmargin=1.5em]
\item ISI Paper~1 --- Technical Specification (v1.0) --- doi:\\
  \url{https://doi.org/10.5281/zenodo.18764227}
\item ISI Paper~2 --- EU-27 Empirical Results 2024 (v1.0) --- doi:\\
  \url{https://doi.org/10.5281/zenodo.18764170}
\end{itemize}
\end{flushleft}
\vfill
\end{titlepage}

\hypersetup{pageanchor=true}
\pagenumbering{arabic}
\pagestyle{fancy}
\fancyhead[L]{\small\sffamily\textit{ISI Country Brief}}
\fancyhead[R]{\small\sffamily\textit{Portugal}}
\renewcommand{\headrulewidth}{0.4pt}

% ---------- Body --------------------------------------------------------------

\section{Executive Summary}
\label{sec:summary}

Portugal ranks~17th in the EU-27 with a composite \ISI score of~0.322, classified as \emph{moderately concentrated}.  The dominant axes are logistics~(0.508) and defence~(0.434).  The structural profile is characterised as dual-axis concentration (logistics and defence).  Under Dirichlet weight perturbation (10\,000~Monte Carlo draws), Portugal's mean rank is~16.95 with a standard deviation of~1.55; the 5th--95th percentile band spans ranks~14 to~19.


\section{Country Position in the EU-27}
\label{sec:position}

Portugal occupies rank~17 of~27 EU member states (composite score~0.322).  The EU-27 mean composite score is~0.344 with a standard deviation of~0.070.  Portugal's score lies below the EU-27 mean, indicating below-average external supplier concentration.  In the ranking, Portugal is situated between Hungary (rank~16; 0.330) and Romania (rank~18; 0.319).


\section{Axis Profile}
\label{sec:axis-profile}

\Cref{tab:axis-profile} presents Portugal's scores on all six \ISI axes.

\begin{table}[htbp]
\centering\small
\caption{Portugal's \ISI axis profile.}
\label{tab:axis-profile}
\begin{tabular}{@{}lr@{}}
\toprule
\textbf{Axis} & \textbf{Score} \\
\midrule
    Finance           & 0.183 \\
    Energy            & 0.349 \\
    Technology        & 0.198 \\
    Defence           & 0.434 \\
    Critical Inputs   & 0.260 \\
    Logistics         & 0.508 \\
\bottomrule
\end{tabular}
\end{table}

The highest axis score is logistics~(0.508), which lies below the EU-27 axis mean of~0.518.  The second-highest axis is defence~(0.434), below the EU-27 axis mean of~0.573.  The lowest axis score is finance~(0.183), above the EU-27 mean of~0.148.


\section{Structural Drivers}
\label{sec:drivers}

\begin{sloppypar}
Portugal's composite score is driven by two axes in combination: logistics~(0.508) and defence~(0.434).  The gap between these two axes is~0.074, indicating a dual-axis concentration profile.  Neither axis alone dominates the composite; rather, the combination of elevated scores on both dimensions determines Portugal's position in the ranking.
\end{sloppypar}


\section{Comparative Context}
\label{sec:comparative}

Across the EU-27, the covariance-based variance decomposition shows that defence (35.1\,\%) and logistics (34.3\,\%) together account for 69.4\,\% of total cross-sectional composite variance.  Portugal's defence score~(0.434) lies below the EU-27 defence mean~(0.573), and its logistics score~(0.508) lies below the EU-27 logistics mean~(0.518).  Under Dirichlet weight perturbation, Portugal's rank standard deviation is~1.55, indicating moderate rank stability.



Portugal's composite of~0.322 lies 0.022 below the EU-27 mean.  Within the axis profile, logistics~(0.508) is the only dimension that exceeds its EU-27 mean, while the remaining five axes score at or below their respective averages.  This pattern---a single above-mean axis paired with broadly below-mean performance elsewhere---is the mirror image of the single-axis-dominant profiles observed at the top of the ranking.

The rank standard deviation of~1.55 and the five-position percentile band~(14--19) reflect moderate stability.  Because logistics is the second-largest contributor to EU-wide composite variance, Portugal's above-mean logistics score provides an upward pull that partially offsets the below-mean defence score.  Under weight schemes that increase the logistics share, Portugal shifts upward; under schemes emphasising defence, it shifts downward.

The composite gap to Hungary~(rank~16; 0.008) is narrow, as is the gap to Romania~(rank~18; 0.003).  Portugal is thus situated in a dense segment of the distribution where ordinal positions are locally interchangeable---a structural feature of the lower-middle ranking rather than a country-specific anomaly.


\section{Interpretation Boundary}
\label{sec:interpretation}

The \ISI measures concentration of external suppliers, not sovereignty in a normative or political sense.  High concentration may reflect geography, economic specialisation, or alliance structures.  A high composite score does not imply policy failure; a low score does not imply policy success.  The index provides an empirical measurement basis --- what policymakers derive from it is a separate question.


% ---------- References --------------------------------------------------------
\clearpage
\vspace*{1.5cm}

\begingroup\emergencystretch=2em\hbadness=10000
\section*{References}

\begin{enumerate}[label={[\arabic*]},leftmargin=2em,itemsep=6pt]
\item International Sovereignty Institute.
  \textit{International Sovereignty Index (ISI) technical specification, version~1.0}.
  Zenodo, 2026.\\ISI Paper Series, Paper~1.
  doi:\,\href{https://doi.org/10.5281/zenodo.18764227}{10.5281/zenodo.18764227}.

\item International Sovereignty Institute.
  \textit{International Sovereignty Index (ISI): EU-27 empirical results 2024 (v1.0)}.
  Zenodo, 2026.\\ISI Paper Series, Paper~2.
  doi:\,\href{https://doi.org/10.5281/zenodo.18764170}{10.5281/zenodo.18764170}.
\end{enumerate}
\endgroup

\end{document}
